\section{Related Work}
\label{sec:related}

\paragraph{Latent world models for planning.}
Learning a compact latent model of the environment and using it for control goes back to World Models~\cite{ha2018world}, PlaNet~\cite{hafner2019planet} and the Dreamer line~\cite{hafner2020dreamer,hafner2021dreamerv2,hafner2025dreamerv3,hafner2025dreamer4}, which train behaviours by actor-critic learning on imagined latent rollouts. Transformer and diffusion world models follow the same recipe in discrete or pixel space~\cite{micheli2023iris,alonso2024diamond}, and JEDI trains a diffusion world model end to end with a JEPA-style latent objective for online reinforcement learning, explicitly without look-ahead planning~\cite{lim2026jedi}. TD-MPC and TD-MPC2 instead plan online with MPPI~\cite{williams2016mppi} in a decoder-free latent space~\cite{hansen2022tdmpc,hansen2024tdmpc2}. Model-based methods with probabilistic ensembles~\cite{chua2018pets,janner2019mbpo,yu2020mopo} are the classical antecedent for treating model uncertainty during planning. Recent surveys cover the wider field~\cite{zidan2026survey,hou2026robotsurvey}.

\paragraph{JEPA world models and their planning objective.}
Joint-embedding predictive architectures predict in representation space rather than pixel space~\cite{assran2023ijepa,bardes2024vjepa}. Action-conditioned variants give zero-shot goal-reaching by planning with a frozen predictor: DINO-WM on pretrained patch features~\cite{zhou2024dinowm}, PLDM on reward-free offline data~\cite{sobal2025pldm}, V-JEPA~2 at video scale~\cite{assran2025vjepa2}, and LeWorldModel (LeWM), which trains an encoder and predictor from pixels with a prediction loss and the SIGReg regulariser of LeJEPA~\cite{maes2026lewm,balestriero2025lejepa}. The design space of this family is mapped in~\cite{terver2025jepawms}, and \texttt{stable-worldmodel} provides shared environments and planners~\cite{maes2026swm}. All of these use a deterministic predictor and rank candidate action sequences by the distance between the predicted terminal latent and the goal latent, usually with CEM~\cite{rubinstein1999cem}.

A fast-growing set of 2026 papers argues that this planning objective is the weak point. The latent distance can correlate poorly with true cost, saturate, or invert at long range~\cite{singh2026objective}, and frozen JEPA objectives can rank candidate actions incorrectly even when closed-loop replanning hides it~\cite{zhang2026arcbench,wang2026decisionmetric,chen2026atm}. Proposed repairs replace or augment the cost with learned temporal or reachability metrics~\cite{li2026trm,bai2026tdjepa,fu2026cgs}, retrieve anchors from offline trajectories~\cite{liu2026anchored}, improve proposals with subgoals~\cite{cheng2026sage}, or learn decision-aligned relations among candidates~\cite{liu2026djepa}. Alrasheed et al.\ study the rollout horizon at which frozen-backbone predictors stop ranking actions reliably~\cite{alrasheed2026limits}. Li et al.\ describe proposal overgeneration: when predicted terminal costs carry residual error, enlarging the candidate pool lowers the feasibility of the minimum-cost candidate~\cite{li2026overgeneration}. This is a selection effect of the kind studied as the optimizer's curse~\cite{smith2006optimizers}, reward-model overoptimization~\cite{gao2023overopt,khalaf2025inferencehacking} and model exploitation~\cite{bhamidipaty2026exploitable}. Our work keeps the distance cost fixed and varies the predictor and how the planner uses it.

\paragraph{Stochastic and distributional world models.}
Stochastic latent states have long been standard in recurrent world models~\cite{hafner2019planet,hafner2025dreamerv3}, whereas the JEPA planners above are deterministic. Several recent papers make the JEPA predictor probabilistic. VJEPA gives a variational formulation, proves that its latent can be a sufficient information state for control, writes the stochastic MPC objective (expected cost under the predictive distribution) and a sampling-based latent MPC routine, and validates on a linear-Gaussian ``noisy TV'' toy~\cite{huang2026vjepa}; it also proves that mean planning is optimal when the predictive covariance does not depend on the action and the cost is quadratic. Var-JEPA gives a related ELBO view~\cite{gogl2026varjepa}. Branch-JEPA represents the transition as a finite set of weighted latent successors trained with the energy score, and is evaluated by offline trajectory forecasting on Argoverse~2 rather than closed-loop control~\cite{song2026branchjepa}. Flow-JEPA replaces LeWM's regression predictor with a flow-matching trajectory generator and plans with CEM, motivated by robustness to visual noise; as described, each candidate is scored with one generated future~\cite{huo2026flowjepa}. FlowWM performs flow matching in pretrained feature space and is evaluated on forecasting and perception benchmarks~\cite{porcher2026flowwm}. Dong proves that observed transitions cannot separate state aliasing from process noise~\cite{dong2026closure}. Training a predictor on a proper scoring rule~\cite{gneiting2007scoring} is standard in probabilistic forecasting: engression fits generative models with the energy score~\cite{shen2025engression}, and AIFS-CRPS trains an operational weather ensemble on the CRPS~\cite{lang2024aifscrps}. Deep ensembles are the other common route to predictive distributions~\cite{lakshminarayanan2017ensembles,chua2018pets}.

\paragraph{Risk-aware planning with learned models.}
PETS propagates particles through a probabilistic ensemble and optimises expected return with MPC~\cite{chua2018pets}; MOPO penalises model uncertainty in offline reinforcement learning~\cite{yu2020mopo}. Risk-sensitive sampling-based MPC uses the conditional value-at-risk~\cite{rockafellar2000cvar} inside MPPI~\cite{yin2022ramppi,enwerem2026riskbelief}, and particle MPC treats any particle representation of uncertainty as a set of scenarios~\cite{dyro2021pmpc}. In latent world models, safety is pursued by conformal error bounds inside robust MPC~\cite{nath2026sls}, by safe reinforcement learning inside world models~\cite{huang2024safedreamer}, and by adaptive latent safety filters~\cite{cao2026worldmodelslie}. D-JEPA includes expected-value and CVaR planning over 16 sampled futures per candidate among its baselines, with nearly identical scores~\cite{liu2026djepa}. Estimating failure probabilities of candidate plans by importance sampling is established in motion planning~\cite{janson2017mcmp,schmerling2017adaptiveis}. For choosing among noisy candidates, the simulation literature offers ranking and selection procedures~\cite{kim2001fully}, racing~\cite{maron1993hoeffding}, and common random numbers~\cite{glasserman1992crn}, and the statistics literature offers empirical-Bayes shrinkage~\cite{efron1973stein}. These tools are standard and we apply them as add-ons to latent planning, so the question we ask is whether they help there.

\paragraph{Test-time adaptation and calibration.}
Adapting a model at test time from a self-supervised signal is well established~\cite{sun2020ttt,wang2021tent}. For latent world models, AdaJEPA updates the predictor inside the MPC loop from observed transitions~\cite{wang2026adajepa}, JEPA-TTT accumulates updates across episodes~\cite{zhang2026jepattt}, and Sandwich-Residuals learns small residual corrections around a frozen predictor~\cite{soni2026sandwich}. All three adapt a deterministic predictor. Adaptive conformal inference~\cite{gibbs2021aci} and conformal decision theory~\cite{lekeufack2024cdt} recalibrate prediction sets or decision parameters online under shift and are the natural tools for recalibrating the spread of a distributional predictor.

\paragraph{Positioning.}
Much of what we use has precedent, and we do not claim it as new. Planning with particles through a probabilistic model is PETS~\cite{chua2018pets}. The stochastic latent MPC objective and a sampling-based planner for JEPA predictors appear in VJEPA~\cite{huang2026vjepa}. Energy-score training of a JEPA predictor appears in Branch-JEPA~\cite{song2026branchjepa}. A generative JEPA predictor planned with CEM on top of LeWM appears in Flow-JEPA~\cite{huo2026flowjepa}. CVaR and expected-value scoring over sampled latent futures appear as baselines in~\cite{liu2026djepa}. The present paper is a controlled study of stochastic latent planning in this setting. We compare deterministic and distributional predictors under the same particle-based planner, on a benchmark with planted stochastic hazards, together with a theoretical account of when planning with a distribution can differ from planning with its mean, and a set of negative results on popular add-ons (risk-averse scoring, shrinkage, racing, common random numbers, importance sampling, and test-time adaptation).

To our knowledge, none of the closed-loop JEPA-planning papers above trains a predictor with a proper scoring rule and evaluates particle-averaged planning against a matched deterministic predictor, and none studies the add-ons listed above systematically. This statement rests on searches of arXiv titles and abstracts and of the web up to 5 October 2026, so it can fail for work we could not find, and it concerns only this combination, not its individual ingredients.
